\documentclass[10pt,a4paper]{amsart}
\usepackage{amsmath,amssymb,amsthm,graphicx}
\usepackage{verbatim}
\newenvironment{reference}{\comment}{\endcomment}
%\newenvironment{reference}{}{}
\newenvironment{slogan}{\comment}{\endcomment}
\newenvironment{history}{\comment}{\endcomment}
% For commutative diagrams we use Xy-pic
\usepackage[all]{xy}

% We use 2cell for 2-commutative diagrams.
\xyoption{2cell}
% Theorem environments.
%
\theoremstyle{plain}
\newtheorem{theorem}[subsection]{Theorem}
\newtheorem{proposition}[subsection]{Proposition}
\newtheorem{lemma}[subsection]{Lemma}

\theoremstyle{definition}
\newtheorem{definition}[subsection]{Definition}
\newtheorem{example}[subsection]{Example}
\newtheorem{exercise}[subsection]{Exercise}
\newtheorem{situation}[subsection]{Situation}

\theoremstyle{remark}
\newtheorem{remark}[subsection]{Remark}
\newtheorem{remarks}[subsection]{Remarks}

\numberwithin{equation}{subsection}

% Macros
%
\def\lim{\mathop{\mathrm{lim}}\nolimits}
\def\colim{\mathop{\mathrm{colim}}\nolimits}
\def\Spec{\mathop{\mathrm{Spec}}}
\def\Hom{\mathop{\mathrm{Hom}}\nolimits}
\def\Ext{\mathop{\mathrm{Ext}}\nolimits}
\def\SheafHom{\mathop{\mathcal{H}\!\mathit{om}}\nolimits}
\def\SheafExt{\mathop{\mathcal{E}\!\mathit{xt}}\nolimits}
\def\Sch{\mathit{Sch}}
\def\Mor{\mathop{\mathrm{Mor}}\nolimits}
\def\Ob{\mathop{\mathrm{Ob}}\nolimits}
\def\Sh{\mathop{\mathit{Sh}}\nolimits}
\def\NL{\mathop{N\!L}\nolimits}
\def\CH{\mathop{\mathrm{CH}}\nolimits}
\def\proetale{{pro\text{-}\acute{e}tale}}
\def\etale{{\acute{e}tale}}
\def\QCoh{\mathit{QCoh}}
\def\Ker{\mathop{\mathrm{Ker}}}
\def\Im{\mathop{\mathrm{Im}}}
\def\Coker{\mathop{\mathrm{Coker}}}
\def\Coim{\mathop{\mathrm{Coim}}}

% Boxtimes
%
\DeclareMathSymbol{\boxtimes}{\mathbin}{AMSa}{"02}

%
% Macros for moduli stacks/spaces
%
\def\QCohstack{\mathcal{QC}\!\mathit{oh}}
\def\Cohstack{\mathcal{C}\!\mathit{oh}}
\def\Spacesstack{\mathcal{S}\!\mathit{paces}}
\def\Quotfunctor{\mathrm{Quot}}
\def\Hilbfunctor{\mathrm{Hilb}}
\def\Curvesstack{\mathcal{C}\!\mathit{urves}}
\def\Polarizedstack{\mathcal{P}\!\mathit{olarized}}
\def\Complexesstack{\mathcal{C}\!\mathit{omplexes}}
% \Pic is the operator that assigns to X its picard group, usage \Pic(X)
% \Picardstack_{X/B} denotes the Picard stack of X over B
% \Picardfunctor_{X/B} denotes the Picard functor of X over B
\def\Pic{\mathop{\mathrm{Pic}}\nolimits}
\def\Picardstack{\mathcal{P}\!\mathit{ic}}
\def\Picardfunctor{\mathrm{Pic}}
\def\Deformationcategory{\mathcal{D}\!\mathit{ef}}
\usepackage{cleveref}
\crefname{originaldefinition}{original definition}{original definitions}
\Crefname{originaldefinition}{Original definition}{Original definitions}
\makeatletter
\newlabel{definition-base}{{definition-base}{source}{}{sourcepointer-1}{}}
\newlabel{definition-base@cref}{{[originaldefinition][0][]definition-base}{[1][0][]source}{}{}{}}
\makeatother
\title{A collection of open sets closed under finite intersections and containing the members of an open cover is a basis}
\author{The Stacks Project Authors}
\date{stacks-project:topology.tex:local-snapshot-20261012; GNU FDL 1.2 or later}
\begin{document}
\maketitle
\noindent\textit{Editorial note.} All delivered mathematical text is the author's own, copied from the pinned Stacks Project chapter source sources/stacks-topology/pinned-original-source.tex; no statement or proof step was rewritten, reordered or removed.\par\medskip
\noindent\textit{Editorial note.} Preamble provenance: the Stacks Project keeps one single shared preamble for every chapter; this card takes it directly from the staged sources/stacks-topology/preamble.tex as three verbatim spans (lines 8-12, 14-18 and 154-217), so no derived file was written and nothing was added or redefined. The wrapper class and the base packages are supplied by the assembler.\par\medskip
\noindent\textit{Editorial note.} Bibliography provenance: neither the counted statement nor its proof contains a citation (this lane rescans the whole span before writing the card), so the card carries no bibliography and no entry was reconstructed.\par\medskip
\noindent\textit{Editorial note.} Selection provenance: this card was opened under the file-based control contract (CONTROL.md), with the mechanical pre-card gate run on its statement and proof interval before the card was written.\par\medskip
\ifdefined\hypertarget\hypertarget{sourcepointer-1}{}\fi
\noindent\textit{Source reference.} Original definition definition-base: sources/stacks-topology/pinned-original-source.tex, source lines 499--499. Navigation pointer only: the counted text refers to the Stacks item definition-base, which is an existing labeled result of this same chapter and is not part of this card. The pointer carries the target's own label as its number and its definition line as its span; its content is not reproduced here.\par\medskip
\begin{lemma}
\label{lemma-make-base}
Let $X$ be a set and let $\mathcal{B}$ be a collection of subsets.
Assume that $X = \bigcup_{B \in \mathcal{B}} B$ and that given
$x \in B_1 \cap B_2$ with $B_1, B_2 \in \mathcal{B}$ there is a
$B_3 \in \mathcal{B}$ with $x \in B_3 \subset B_1 \cap B_2$.
Then there is a unique topology on $X$ such that $\mathcal{B}$
is a basis for this topology.
\end{lemma}

\begin{proof}
Let $\sigma(\mathcal B)$ be the set of subsets of $X$ which can be written
as unions of elements of $\mathcal B$. We claim $\sigma(\mathcal{B})$
is a topology. Namely, the empty set is an element of $\sigma(\mathcal{B})$
(as an empty union) and $X$ is an element of $\sigma(\mathcal{B})$ (as the
union of all elements of $\mathcal{B}$). It is clear that $\sigma(\mathcal{B})$
is preserved under unions. Finally, if $U, V \in \sigma(\mathcal{B})$
then write $U = \bigcup_{i \in I} U_i$ and $V = \bigcup_{j \in J} V_j$
with $U_i, V_j \in \mathcal{B}$ for all $i \in I$ and $j \in J$.
Then
$$
U \cap V = \bigcup\nolimits_{i \in I, j \in J}
U_i \cap V_j
$$
The assumption in the lemma tells us that $U_i \cap V_j \in \sigma(\mathcal{B})$
hence we see that $U \cap V$ is too. Thus $\sigma(\mathcal{B})$ is a topology.
Properties (1) and (2) of Definition \ref{definition-base} are
immediate for this topology. To prove the uniqueness of this topology
let $\mathcal T$ be a topology on $X$ such that $\mathcal B$ is a base for it.
Then of course every element of $\mathcal{B}$ is in $\mathcal{T}$ by
(1) of Definition \ref{definition-base} and
hence $\sigma(\mathcal{B}) \subset \mathcal{T}$. Conversely, part
(2) of Definition \ref{definition-base} tells us that every element
of $\mathcal{T}$ is a union of elements of $\mathcal{B}$, i.e.,
$\mathcal{T} \subset \sigma(\mathcal{B})$. This finishes the proof.
\end{proof}

\end{document}
